\documentclass[11pt,a4paper,fontset=fandol]{ctexart}

\usepackage[a4paper,top=2.25cm,bottom=2.45cm,left=2.25cm,right=2.25cm,headheight=18pt,headsep=0.55cm,footskip=24pt]{geometry}
\usepackage{amsmath,amssymb,amsthm}
\usepackage{fancyhdr}
\usepackage{lastpage}
\usepackage{enumitem}
\usepackage{titlesec}
\usepackage{xcolor}
\usepackage{hyperref}
\usepackage{bookmark}
\usepackage[most]{tcolorbox}
\usepackage{setspace}
\usepackage{array}

\hypersetup{
  colorlinks=true,
  linkcolor=blue!50!black,
  urlcolor=blue!50!black
}

\setstretch{1.13}
\setlength{\parindent}{2em}
\setlength{\parskip}{0.25em}
\allowdisplaybreaks
\setlist[enumerate]{itemsep=0.45em, topsep=0.35em, leftmargin=2.4em}
\setlist[itemize]{itemsep=0.25em, topsep=0.25em, leftmargin=1.9em}

\definecolor{mainblue}{RGB}{36,83,140}
\definecolor{lightblue}{RGB}{241,246,252}
\definecolor{lightgray}{RGB}{246,246,246}
\definecolor{rulegray}{RGB}{110,118,128}

\titleformat{\section}
  {\Large\bfseries\color{mainblue}}
  {\thesection.}
  {0.45em}
  {}
  [\vspace{0.2em}\color{rulegray}\titlerule]
\titlespacing*{\section}{0pt}{1.1em}{0.7em}

\pagestyle{fancy}
\fancyhf{}
\fancyhead[L]{\small 函数图像平移周测 B 卷}
\fancyhead[C]{\small 代数专题：综合判断与变式}
\fancyhead[R]{\small 刘文博}
\fancyfoot[C]{\small 第 \thepage 页 / 共 \pageref{LastPage} 页}
\renewcommand{\headrulewidth}{0.55pt}
\renewcommand{\footrulewidth}{0.35pt}

\newtcolorbox{infobox}[1]{
  breakable,
  colback=lightblue,
  colframe=mainblue,
  boxrule=0.7pt,
  arc=1mm,
  title=\textbf{#1},
  fonttitle=\bfseries
}

\newenvironment{solution}{\par\noindent\textbf{参考解答：}}{\hfill$\square$\par}
\newcommand{\answerspace}{\vspace{2.3cm}}
\newcommand{\biganswerspace}{\vspace{3.2cm}}

\begin{document}

\thispagestyle{empty}
\begin{center}
{\fontsize{22}{28}\selectfont\bfseries 函数图像平移周测 B 卷\par}
\vspace{0.4cm}
{\large 适合初中阶段函数专题提高训练\par}
\end{center}

\begin{infobox}{考试说明}
\begin{itemize}
  \item 测试时间：70--75 分钟
  \item 满分：100 分
  \item B 卷更强调方法迁移，请尽量说明“为什么这样平移”，不要只写最后结果。
\end{itemize}
\end{infobox}

\section{试题}

\begin{enumerate}
  \item （12 分）请用“图像上的点跟着移动”的方法说明：为什么函数
  \[
  y=f(x)
  \]
  向右平移 $a$ 个单位后，新方程应写成
  \[
  y=f(x-a).
  \]
  \biganswerspace

  \item （12 分）把函数
  \[
  y=x^2-2x
  \]
  向右平移 $3$ 个单位，再向下平移 $2$ 个单位，求化简后的新方程。
  \answerspace

  \item （12 分）已知函数
  \[
  y=|x|
  \]
  经过平移后，顶点变成 $(4,-3)$。求新方程。
  \answerspace

  \item （12 分）已知函数
  \[
  y=f(x)
  \]
  经过平移后变成
  \[
  y=f(x+2)-5.
  \]
  如果原图像上有点 $(3,7)$，那么平移后对应点的坐标是什么？
  \answerspace

  \item （12 分）已知函数
  \[
  y=(x-4)^2+6,
  \]
  说出它是由
  \[
  y=x^2
  \]
  怎样平移得到的。
  \answerspace

  \item （12 分）已知函数
  \[
  y=x^2-8x+13,
  \]
  说出它是由
  \[
  y=x^2
  \]
  怎样平移得到的。
  \answerspace

  \item （14 分）将直线
  \[
  y=2x+1
  \]
  先向右平移 $1$ 个单位，再作一次竖直平移后，恰好经过原点。求第二次平移的方向、距离以及最终方程。
  \biganswerspace

  \item （14 分）已知原图像上的点 $(1,2)$ 在函数
  \[
  y=f(x)
  \]
  上，平移后该点对应到新图像上的点 $(5,-1)$。求新图像的方程，并说明平移方向和距离。
  \biganswerspace
\end{enumerate}

\newpage
\section{详细解答}

\begin{enumerate}
  \item \begin{solution}
  设原图像上有一点 $(x_0,y_0)$，因为它在函数
  \[
  y=f(x)
  \]
  上，所以有
  \[
  y_0=f(x_0).
  \]

  如果整张图像向右平移 $a$ 个单位，那么这个点会变成
  \[
  (x_0+a,y_0).
  \]

  现在把新图像上的横坐标记作 $x$，那么它对应原图像上的横坐标应是
  \[
  x-a.
  \]
  因此新图像上的函数值应为
  \[
  y=f(x-a).
  \]

  所以，函数 $y=f(x)$ 向右平移 $a$ 个单位后，新方程应写成
  \[
  \boxed{y=f(x-a)}.
  \]
  \end{solution}

  \item \begin{solution}
  先向右平移 $3$ 个单位，把整个式子里的 $x$ 换成 $x-3$：
  \[
  y=(x-3)^2-2(x-3).
  \]
  再向下平移 $2$ 个单位：
  \[
  y=(x-3)^2-2(x-3)-2.
  \]
  展开化简：
  \[
  y=x^2-6x+9-2x+6-2=x^2-8x+13.
  \]
  所以新方程为
  \[
  \boxed{y=x^2-8x+13}.
  \]
  \end{solution}

  \item \begin{solution}
  函数 $y=|x|$ 的顶点原来是 $(0,0)$。

  现在顶点变成 $(4,-3)$，说明图像向右平移了 $4$ 个单位，再向下平移了 $3$ 个单位。

  因此新方程为
  \[
  \boxed{y=|x-4|-3}.
  \]
  \end{solution}

  \item \begin{solution}
  方程
  \[
  y=f(x+2)-5
  \]
  表示图像向左平移 $2$ 个单位，再向下平移 $5$ 个单位。

  所以点 $(3,7)$ 平移后，横坐标减 $2$，纵坐标减 $5$，得到
  \[
  \boxed{(1,2)}.
  \]
  \end{solution}

  \item \begin{solution}
  方程
  \[
  y=(x-4)^2+6
  \]
  已经直接告诉我们：从 $y=x^2$ 出发，图像先向右平移 $4$ 个单位，再向上平移 $6$ 个单位。
  \end{solution}

  \item \begin{solution}
  先配方：
  \[
  y=x^2-8x+13=(x-4)^2-3.
  \]
  因此它是由 $y=x^2$ \textbf{向右平移 $4$ 个单位，再向下平移 $3$ 个单位}得到的。
  \end{solution}

  \item \begin{solution}
  第一次向右平移 $1$ 个单位后，直线变为
  \[
  y=2(x-1)+1=2x-1.
  \]

  现在它还没有经过原点，因为当 $x=0$ 时，
  \[
  y=-1.
  \]

  要让它经过原点，就需要再向上平移 $1$ 个单位。

  所以第二次平移是\textbf{向上平移 $1$ 个单位}，最终方程为
  \[
  y=2x-1+1=2x.
  \]
  即
  \[
  \boxed{y=2x}.
  \]
  \end{solution}

  \item \begin{solution}
  点 $(1,2)$ 平移后变成 $(5,-1)$。

  横坐标从 $1$ 变到 $5$，增加了 $4$，说明向右平移 $4$ 个单位；

  纵坐标从 $2$ 变到 $-1$，减少了 $3$，说明向下平移了 $3$ 个单位。

  因此新图像方程为
  \[
  \boxed{y=f(x-4)-3}.
  \]

  所以平移方式是：\textbf{向右平移 $4$ 个单位，再向下平移 $3$ 个单位}。
  \end{solution}
\end{enumerate}

\end{document}
